\section{Giới thiệu}

\begin{frame}
\frametitle{Bối cảnh}
\framesubtitle{Giám sát giao thông thông minh}

\begin{itemize}[<+->]
  \item Đô thị hóa và mật độ phương tiện tăng nhanh gây áp lực lên hệ thống
        giao thông
  \item Camera giám sát tạo ra nguồn dữ liệu video phong phú, chi phí thấp và
        dễ triển khai hơn nhiều cảm biến chuyên dụng
  \item Phát hiện và đếm phương tiện hỗ trợ giám sát lưu lượng, phát hiện ùn
        tắc và quy hoạch hạ tầng giao thông
  \item Học sâu giúp các hệ thống thị giác máy tính đạt độ chính xác cao hơn
        trong môi trường thực tế
\end{itemize}

\vspace{0.2cm}

{\scriptsize Nguồn: \cite{YANG2018143,song2019}}

\end{frame}

\begin{frame}
\frametitle{Phát biểu bài toán}
\framesubtitle{Vehicle Detection and Counting}

\begin{itemize}[<+->]
  \item Đầu vào: video hoặc chuỗi ảnh từ camera giám sát giao thông cố định
  \item Phát hiện vị trí phương tiện trong từng khung hình bằng hộp giới hạn
  \item Theo dõi phương tiện qua các khung hình để duy trì định danh
  \item Đếm số lượng phương tiện đi qua vùng quan sát xác định trước
  \item Đầu ra: vị trí, quỹ đạo và số lượng phương tiện trong khoảng thời gian
        quan sát
\end{itemize}

\end{frame}

\begin{frame}
\frametitle{Thách thức}
\framesubtitle{Môi trường giao thông thực tế}

\begin{itemize}[<+->]
  \item Điều kiện ánh sáng và thời tiết thay đổi theo thời gian
  \item Che khuất giữa các phương tiện trong cảnh giao thông đông
  \item Góc nhìn camera, phối cảnh và kích thước phương tiện không đồng nhất
  \item Nhiều chủng loại phương tiện làm tăng độ khó cho phát hiện và theo dõi
  \item Hệ thống cần cân bằng giữa độ chính xác và khả năng xử lý gần thời
        gian thực
\end{itemize}

\vspace{0.2cm}

{\scriptsize Nguồn: \cite{YANG2018143,dai2019}}

\end{frame}

\begin{frame}
\frametitle{Các hướng tiếp cận liên quan}
\framesubtitle{Từ thị giác cổ điển đến học sâu}

\begin{itemize}[<+->]
  \item Phương pháp cổ điển: chênh lệch khung hình, luồng quang học và trừ nền
  \item Đặc trưng thủ công như SIFT, HOG, Haar-like kết hợp bộ phân loại SVM
        hoặc Adaboost
  \item Mô hình hai giai đoạn như R-CNN và Faster R-CNN ưu tiên độ chính xác
        nhưng khó đáp ứng thời gian thực
  \item Mô hình một giai đoạn như YOLO và SSD cân bằng tốt hơn giữa tốc độ và
        độ chính xác
\end{itemize}

\vspace{0.2cm}

{\scriptsize Nguồn: \cite{song2019,dai2019,yaseen2024yolov8}}

\end{frame}

\begin{frame}
\frametitle{Hướng giải quyết của đồ án}
\framesubtitle{Detection-based tracking and counting}

\begin{itemize}[<+->]
  \item Fine-tune YOLOv8n trên UA-DETRAC để phát hiện phương tiện
  \item Sử dụng tracker kiểu detection-based để ghép nối phương tiện theo thời
        gian
  \item So sánh SORT và BYTE, chọn cấu hình phù hợp cho tracking
  \item Đếm phương tiện dựa trên quỹ đạo sau tracking và vùng quan sát
  \item Đánh giá thực nghiệm theo từng thành phần: detection, tracking và
        counting
\end{itemize}

\vspace{0.2cm}

{\scriptsize Nguồn dữ liệu: \cite{wen2020uadetrac}}

\end{frame}
